% Options for packages loaded elsewhere
% Options for packages loaded elsewhere
\PassOptionsToPackage{unicode}{hyperref}
\PassOptionsToPackage{hyphens}{url}
\PassOptionsToPackage{dvipsnames,svgnames,x11names}{xcolor}
%
\documentclass[
  english,
  11pt,
]{article}
\usepackage{xcolor}
\usepackage[margin=25mm]{geometry}
\usepackage{amsmath,amssymb}
\setcounter{secnumdepth}{-\maxdimen} % remove section numbering
\usepackage{iftex}
\ifPDFTeX
  \usepackage[T1]{fontenc}
  \usepackage[utf8]{inputenc}
  \usepackage{textcomp} % provide euro and other symbols
\else % if luatex or xetex
  \usepackage{unicode-math} % this also loads fontspec
  \defaultfontfeatures{Scale=MatchLowercase}
  \defaultfontfeatures[\rmfamily]{Ligatures=TeX,Scale=1}
\fi
\usepackage{lmodern}
\ifPDFTeX\else
  % xetex/luatex font selection
\fi
% Use upquote if available, for straight quotes in verbatim environments
\IfFileExists{upquote.sty}{\usepackage{upquote}}{}
\IfFileExists{microtype.sty}{% use microtype if available
  \usepackage[]{microtype}
  \UseMicrotypeSet[protrusion]{basicmath} % disable protrusion for tt fonts
}{}
\makeatletter
\@ifundefined{KOMAClassName}{% if non-KOMA class
  \IfFileExists{parskip.sty}{%
    \usepackage{parskip}
  }{% else
    \setlength{\parindent}{0pt}
    \setlength{\parskip}{6pt plus 2pt minus 1pt}}
}{% if KOMA class
  \KOMAoptions{parskip=half}}
\makeatother
% Make \paragraph and \subparagraph free-standing
\makeatletter
\ifx\paragraph\undefined\else
  \let\oldparagraph\paragraph
  \renewcommand{\paragraph}{
    \@ifstar
      \xxxParagraphStar
      \xxxParagraphNoStar
  }
  \newcommand{\xxxParagraphStar}[1]{\oldparagraph*{#1}\mbox{}}
  \newcommand{\xxxParagraphNoStar}[1]{\oldparagraph{#1}\mbox{}}
\fi
\ifx\subparagraph\undefined\else
  \let\oldsubparagraph\subparagraph
  \renewcommand{\subparagraph}{
    \@ifstar
      \xxxSubParagraphStar
      \xxxSubParagraphNoStar
  }
  \newcommand{\xxxSubParagraphStar}[1]{\oldsubparagraph*{#1}\mbox{}}
  \newcommand{\xxxSubParagraphNoStar}[1]{\oldsubparagraph{#1}\mbox{}}
\fi
\makeatother


\usepackage{longtable,booktabs,array}
\newcounter{none} % for unnumbered tables
\usepackage{calc} % for calculating minipage widths
% Correct order of tables after \paragraph or \subparagraph
\usepackage{etoolbox}
\makeatletter
\patchcmd\longtable{\par}{\if@noskipsec\mbox{}\fi\par}{}{}
\makeatother
% Allow footnotes in longtable head/foot
\IfFileExists{footnotehyper.sty}{\usepackage{footnotehyper}}{\usepackage{footnote}}
\makesavenoteenv{longtable}
\usepackage{graphicx}
\makeatletter
\newsavebox\pandoc@box
\newcommand*\pandocbounded[1]{% scales image to fit in text height/width
  \sbox\pandoc@box{#1}%
  \Gscale@div\@tempa{\textheight}{\dimexpr\ht\pandoc@box+\dp\pandoc@box\relax}%
  \Gscale@div\@tempb{\linewidth}{\wd\pandoc@box}%
  \ifdim\@tempb\p@<\@tempa\p@\let\@tempa\@tempb\fi% select the smaller of both
  \ifdim\@tempa\p@<\p@\scalebox{\@tempa}{\usebox\pandoc@box}%
  \else\usebox{\pandoc@box}%
  \fi%
}
% Set default figure placement to htbp
\def\fps@figure{htbp}
\makeatother



\ifLuaTeX
\usepackage[bidi=basic,shorthands=off]{babel}
\else
\usepackage[bidi=default,shorthands=off]{babel}
\fi
\ifLuaTeX
  \usepackage{selnolig} % disable illegal ligatures
\fi


\setlength{\emergencystretch}{3em} % prevent overfull lines

\providecommand{\tightlist}{%
  \setlength{\itemsep}{0pt}\setlength{\parskip}{0pt}}



 


\makeatletter
\@ifpackageloaded{caption}{}{\usepackage{caption}}
\AtBeginDocument{%
\ifdefined\contentsname
  \renewcommand*\contentsname{Table of contents}
\else
  \newcommand\contentsname{Table of contents}
\fi
\ifdefined\listfigurename
  \renewcommand*\listfigurename{List of Figures}
\else
  \newcommand\listfigurename{List of Figures}
\fi
\ifdefined\listtablename
  \renewcommand*\listtablename{List of Tables}
\else
  \newcommand\listtablename{List of Tables}
\fi
\ifdefined\figurename
  \renewcommand*\figurename{Figure}
\else
  \newcommand\figurename{Figure}
\fi
\ifdefined\tablename
  \renewcommand*\tablename{Table}
\else
  \newcommand\tablename{Table}
\fi
}
\@ifpackageloaded{float}{}{\usepackage{float}}
\floatstyle{ruled}
\@ifundefined{c@chapter}{\newfloat{codelisting}{h}{lop}}{\newfloat{codelisting}{h}{lop}[chapter]}
\floatname{codelisting}{Listing}
\newcommand*\listoflistings{\listof{codelisting}{List of Listings}}
\makeatother
\makeatletter
\makeatother
\makeatletter
\@ifpackageloaded{caption}{}{\usepackage{caption}}
\@ifpackageloaded{subcaption}{}{\usepackage{subcaption}}
\makeatother
\usepackage{bookmark}
\IfFileExists{xurl.sty}{\usepackage{xurl}}{} % add URL line breaks if available
\urlstyle{same}
% fallback for those not using the hyperref driver hyperxmp:
\makeatletter
\@ifundefined{xmpquote}{\newcommand{\xmpquote}[1]{#1}}{}
\makeatother
\hypersetup{
  pdftitle={Convex integration: how tiny ripples can solve a big constraint},
  pdflang={en},
  colorlinks=true,
  linkcolor={blue},
  filecolor={Maroon},
  citecolor={Blue},
  urlcolor={Blue},
  pdfcreator={LaTeX via pandoc}}


\title{Convex integration: how tiny ripples can solve a big constraint}
\usepackage{etoolbox}
\makeatletter
\providecommand{\subtitle}[1]{% add subtitle to \maketitle
  \apptocmd{\@title}{\par {\large #1 \par}}{}{}
}
\makeatother
\subtitle{From a zigzag line to shear, averages, and the paper's
oscillations}
\author{}
\date{2026-09-12}
\begin{document}
\maketitle


\href{../paper-map/article.html}{Start with the paper map} · Main route,
article 5 of 9

Suppose a shape is nearly what you want, but it fails a local rule.
Changing the whole shape might ruin everything already built. Could
adding very fine detail fix the local rule while keeping the overall
shape almost unchanged?

That question leads to a central idea of \textbf{convex integration}. It
is a way of constructing objects that satisfy difficult local
constraints by using carefully arranged fine variations. In fluid
problems, the variations can also create a useful averaged momentum
transport.

We will begin with slopes and averages, then connect them to two places
in the OpenAI paper: preparing the background shear in Appendix C and
realizing the missing stress with pulses in Section 7. They share a
strategy, but the paper must also make the pulses obey their viscous
dynamics.

\subsection{A line that looks flat but keeps climbing and
descending}\label{a-line-that-looks-flat-but-keeps-climbing-and-descending}

Imagine drawing from the left side of a page to the right while
following a peculiar rule: every straight segment must have slope \(+1\)
or \(-1\).

A horizontal segment is forbidden. Yet you can go up for a short
distance, then down for the same distance, returning to your original
height. Repeat this many times. From far away, the zigzag looks almost
horizontal.

Slope is just ``rise divided by run.'' If an upward segment runs right
by 0.01 units, it rises by 0.01 units. Its downward partner returns to
the original height. Make the segments ten times shorter and the
zigzag's height becomes ten times smaller. The slopes remain \(+1\) and
\(-1\).

\textbf{The height can be tiny while the local slopes remain
substantial.}

\begin{figure}

\centering{

\pandocbounded{\includegraphics[keepaspectratio,alt={Coarse and fine triangular zigzags approach a horizontal line; a separate average diagram shows how plus one and minus one give zero.}]{figures/reader.png}}

}

\caption{\label{fig-ripples}Successively finer zigzags stay closer to
the horizontal line while retaining slopes of plus or minus one between
corners. This is a slope analogy, not a fluid velocity field.}

\end{figure}%

The corners need a qualification: the slope is undefined right at each
sharp turn. This elementary picture is a piecewise-linear example.
Smooth fluid constructions need more carefully designed families and
transitions. The picture teaches the separation between a small change
in values and a large change in slopes.

\subsection{Why the word ``convex''
appears}\label{why-the-word-convex-appears}

Spend half the horizontal distance going up and half going down. The
average slope is

\[
\frac12(+1)+\frac12(-1)=0.
\]

Spend three quarters going up and one quarter going down, and the
average becomes \(0.5\). By changing the proportions, we can obtain any
average between \(-1\) and \(+1\).

The permitted local slopes are the two endpoints. Their \emph{convex
hull} is the entire interval between them: all the weighted averages
obtained with nonnegative weights that add to one.

This exposes the useful gap. A coarse description might allow an average
of zero even though zero is not an allowed local slope. Fine structure
can make the local slopes obey their stricter rule while preserving the
coarse average.

For this example, ``integration'' can be understood as building the
height by adding up the little rises and falls. The mathematical method
extends this idea to much richer constraints. This is an entry point,
not its full definition. \href{https://arxiv.org/abs/math/0702079}{De
Lellis and Székelyhidi, fluid equations as a differential inclusion}

\subsection{Where that idea appears directly in the
paper}\label{where-that-idea-appears-directly-in-the-paper}

The background vortex must have suitable \textbf{shear}: neighbouring
radii must move differently in a way that supports the required pulses.
The original joined profile satisfies a weaker set of conditions in part
of the annulus, but the waves need stronger local conditions.

Appendix C constructs a smoothly varying periodic family of acceptable
shear values. Their average equals the original shear. It then inserts
rapid radial variation and adds up its effect to obtain the modified
velocity profiles.

The pattern is the one we just met:

\begin{itemize}
\tightlist
\item
  The average supplies the large-scale target.
\item
  Fine variation supplies the stronger local behaviour.
\item
  Adding up the variations makes only a small change to the profile
  values.
\item
  A separate repair restores the required weighted totals.
\end{itemize}

The actual constraints involve two coupled shear components, not the two
slopes in our sketch. The Appendix C construction also uses a
sufficiently large \textbf{finite} frequency; it produces a smooth
profile. Its role is to make the background compatible with the later
waves.
\href{https://cdn.openai.com/pdf/32d9f210-8b73-45e0-91bc-82a30aef8a9a/navier-stokes.pdf\#page=158}{Paper,
Appendix C.1--C.2, especially (C.1) and (C.12)}

This is a particularly direct way to see the convex-integration idea in
the paper. It deserves its own explanation, rather than being hidden
behind a phrase such as ``admissible stress cone.''

\subsection{What a ``stress cone'' is trying to tell
us}\label{what-a-stress-cone-is-trying-to-tell-us}

Think of two ingredients that can be mixed only in nonnegative amounts.
Each contributes a particular pair of effects. Their mixtures fill a set
of achievable pairs.

That is the elementary meaning of the relevant \emph{cone}: the
collection of transport combinations the chosen pulse families can
supply. The desired transport must lie inside it. You cannot demand a
combination that the available ingredients cannot produce.

The restriction comes from physics and algebra. The weight of a wave's
leading quadratic contribution is proportional to its amplitude
\textbf{squared}, which is nonnegative. The transport directions,
however, may contain signed components. Positivity of a weight does not
mean all forces point the same way.

Appendix C prepares the background so that its required transport fits
the available pulse families. Section 7 then chooses their weights. The
background and the waves must be designed together.

\subsection{An average velocity can miss a real
effect}\label{an-average-velocity-can-miss-a-real-effect}

Take two equal groups with velocity fluctuations \(+2\) and \(-2\) in
one chosen direction. Their mean fluctuation is zero. Their squared
fluctuations are both 4, so the mean square is 4.

Thus these two operations give different answers:

\[
(\text{average fluctuation})^2=0,
\qquad
\text{average of fluctuation squared}=4.
\]

Fluid momentum transport contains products of velocity components. As a
result, fine motion may disappear from the averaged velocity while still
contributing to averaged transport.

This is the bridge from fine structure to fluid stress. It is not enough
to say ``the ripples are small, so they do nothing.'' Their size,
spacing, direction, and the quantities being averaged all matter. The
\href{../d4_mean_stress/article.html}{next article} works through an
actual transport example using two signs and multiplication.

In a fluid convex-integration construction, one commonly begins with an
approximate flow and an extra stress describing what remains to be
supplied. Fine velocity corrections account for that stress while
creating smaller errors to be handled later.

\subsection{Why adding arbitrary ripples does not solve
Navier--Stokes}\label{why-adding-arbitrary-ripples-does-not-solve-navierstokes}

The ripples must themselves satisfy fluid constraints. They cannot
create or destroy fluid volume, violate the required boundary behaviour,
or leave uncontrolled forces behind.

Viscosity makes very short waves especially costly: fine velocity
differences are smoothed rapidly. The OpenAI construction therefore
includes a specific pulse life cycle. Background shear amplifies a tiny
seed; later the changing wave pattern strengthens viscous damping and
the pulse fades.

It must also account for the force used to seed the pulses, the small
errors introduced when turning them on and off, and interactions with
other corrections. Matching an average stress is a vital step, but it
leaves all of this work to do.
\href{https://cdn.openai.com/pdf/32d9f210-8b73-45e0-91bc-82a30aef8a9a/navier-stokes.pdf}{Paper,
Sections 2.2, 6--9}

\subsection{The difficult part is controlling what happens in the
limit}\label{the-difficult-part-is-controlling-what-happens-in-the-limit}

Imagine reducing a mismatch from 1 to \(1/2\), then \(1/4\), then
\(1/8\). That sequence suggests a route to zero. A real construction
must prove more: that the corrections combine into a field with the
required properties, and that their interactions do not reintroduce a
large error.

Fine ripples reveal the danger. Their heights can shrink while their
slopes stay large. Their slopes may even change ever more rapidly. Small
velocity corrections therefore do not automatically mean small
acceleration or viscous terms.

Some classical fluid convex-integration results construct \textbf{weak
solutions}: solutions understood through balances averaged against
smooth probes, without requiring a classical velocity derivative at
every point. Buckmaster and Vicol's Navier--Stokes nonuniqueness result
belongs to that setting. A weak solution is a mathematical notion, not a
synonym for a poorly resolved numerical simulation.
\href{https://arxiv.org/abs/1709.10033}{Buckmaster and Vicol}

The target of the paper we are explaining is different: a flow smooth
before the singular time and a force that remains smooth through it. Its
derivative and summation estimates must establish that stronger
time-dependent structure. Calling something ``convex integration'' does
not remove those obligations.

\subsection{Where this fits in the bigger
picture}\label{where-this-fits-in-the-bigger-picture}

There are two related uses of fine structure to remember.
\textbf{Appendix C changes local shear while preserving the large-scale
profile and repairing its totals. The later pulses create the averaged
momentum transport needed to repair the background force balance.}

The first prepares suitable conditions for the second. Growth, viscous
decay, and repeated correction then make the oscillatory construction
work as part of the paper's full argument.

Continue with \href{../d4_mean_stress/article.html}{zero average,
nonzero momentum transport}, or return to the
\href{../paper-map/article.html}{paper map}. The
\href{source-notes.md}{source notes} distinguish the elementary
illustrations, the general convex-integration viewpoint, and the
source's specific construction.




\end{document}
